\documentclass[10pt,a4paper]{amsart}
\usepackage[margin=19mm]{geometry}
\usepackage{amsmath,amssymb,mathtools}
\usepackage[scr=rsfs,cal=euler]{mathalfa}
\usepackage{xfrac}
\usepackage[unicode,hidelinks,bookmarks=false]{hyperref}
\newcommand{\ie}{i.e.\ }
\newcommand{\cf}{cf.\ }
\newcommand{\resp}{respectively\ }
\newcommand{\Subsection}{\subsection}
\providecommand{\nobreakdash}{\nobreak\hbox{-}}
\setlength{\emergencystretch}{2em}
\multlinegap0pt
\newcommand{\C}{\mathcal C}
\newcommand{\Hom}{\mathrm{Hom}}
\newcommand{\Loc}{\mathrm{Loc}}
\newcommand{\SLoc}{\operatorname{SLoc}}
\newcommand{\id}{\mathrm{id}}
\newcommand{\one}{1\hskip-3.5pt1}
\theoremstyle{plain}
\newtheorem{theorem}{Theorem}
\newtheorem{proposition}[theorem]{Proposition}
\newtheorem{lemma}[theorem]{Lemma}
\newtheorem{corollary}[theorem]{Corollary}
\theoremstyle{definition}
\newtheorem{definition}[theorem]{Definition}
\newtheorem{assumption}[theorem]{Assumption}
\theoremstyle{remark}
\newtheorem{remark}[theorem]{Remark}
\title[A categorical and algebro-geometric theory of localization]{A categorical and algebro-geometric theory of localization — Proposition 4.13}
\author{Mauricio Corr\^ea, Simone Noja}
\date{Source: arXiv:2604.03845v5 (CC BY 4.0)}
\begin{document}
\maketitle

\noindent\emph{Source and licence.} The delivered work is the article shown in the title above, version arXiv:2604.03845v5 (CC BY 4.0), distributed under the Creative Commons Attribution 4.0 International licence, as recorded in the arXiv licence metadata for this exact version and shown on the abstract page saved with this item. The full licence notice, which carries the licence URL and the address of that abstract page, is the bundled file \texttt{licenses/arxiv-2604.03845-CC-BY-4.0.txt}.\par\smallskip

\noindent\textit{Editorial scope.} Counted unit: the article's Proposition 4.13 (source0.tex lines 653--667, source label \texttt{prop:Loc-excision}); statement and proof are the article's own text line for line, with no line omitted, substituted or re-flowed. Required context retained uncounted: the article's recollement definition (source0.tex lines 336--349), which the counted statement invokes. Editorial formatting only: an AMS \texttt{amsart} preamble that restates the article's own macro definitions (\texttt{\textbackslash C}, \texttt{\textbackslash Hom}, \texttt{\textbackslash Loc}, \texttt{\textbackslash SLoc}, \texttt{\textbackslash id}, \texttt{\textbackslash one}) and the theorem environments the retained lines use; the article ships no class or style file, so none is shipped or input; sequential numbering of the retained environments in order of appearance, whereas the article prints them with its own section-relative numbers; equation numbering is not reproduced and every cross-reference inside the retained text resolves to the wrapper's numbering. No citation occurs inside any retained line, so the wrapper carries no bibliography. No asset-inclusion command, no drawing environment and no image asset occurs in this package.


\section*{Retained context: def:recollement (source0.tex lines 336--349)}

\begin{definition} \label{def:recollement}
We assume:
\begin{itemize}
\item adjunctions $i^*\dashv i_*\dashv i^!$ and $j_!\dashv j^*\dashv j_*$;
\item full faithfulness of $i_*$, $j_!$ and $j_*$;
\item vanishings $j^*i_*=0$, $i^*j_!=0$ and $i^!j_*=0$;
\item functorial distinguished triangles, for each $M\in\C(X)$:
$$
j_!j^*M \longrightarrow M \longrightarrow i_*i^*M \xrightarrow{+1},
\qquad
i_*i^!M \longrightarrow M \longrightarrow j_*j^*M \xrightarrow{+1}.
$$
\end{itemize}
\end{definition}


\section{The counted unit: Proposition 4.13 and its complete original proof}

\begin{proposition} \label{prop:Loc-excision}
Assume the recollement hypotheses of Definition~\ref{def:recollement}. Let $v:V\hookrightarrow X$ be an open neighbourhood of $Z$ and write
$i_V:Z\hookrightarrow V$ for the induced closed immersion and $j_V:V\setminus Z\hookrightarrow V$ for its open complement.
Assume (BC) for the square determined by $v$ and $i$ (so that $v^*i_*\simeq i_{V*}$ and $i_V^!v^*\simeq v_Z^*i^!$).
Let $A\in\C(X)$, fix $d\in\mathbb Z$, and let $c\in H^d(X;A)$ satisfy $j^*c=0$.
Then $j_V^*(v^*c)=0$ and, under the canonical identification
$$
\Hom_{\C(Z)}(\one_Z,i^!A[d])\ \simeq\ \Hom_{\C(Z)}(\one_Z,i_V^!v^*A[d]),
$$
one has an equality of torsors
$$
\Loc_Z^{\mathrm{tor},X}(c)\ =\ \Loc_Z^{\mathrm{tor},V}(v^*c).
$$
In particular, if either side is a singleton, then so is the other and $\Loc_Z^{X}(c)=\Loc_Z^{V}(v^*c)$.
\end{proposition}

\begin{proof}
The vanishing is immediate from functoriality:
$$
j_V^*(v^*c)= (j^*c)|_{V\setminus Z}=0.
$$

For cohomology with supports, adjunction identifies
\begin{align}
H_Z^d(X;A)
&=\Hom_{\C(X)}(\one_X,i_*i^!A[d])
\cong \Hom_{\C(Z)}(i^*\one_X,i^!A[d])
\cong \Hom_{\C(Z)}(\one_Z,i^!A[d]), \label{eq:excision-adj-X}\\
H_Z^d(V;v^*A)
&=\Hom_{\C(V)}(\one_V,i_{V*}i_V^!v^*A[d])
\cong \Hom_{\C(Z)}(i_V^*\one_V,i_V^!v^*A[d])
\cong \Hom_{\C(Z)}(\one_Z,i_V^!v^*A[d]). \label{eq:excision-adj-V}
\end{align}
By Beck--Chevalley for the square determined by $v$ and $i$, one has a canonical isomorphism
\begin{equation}\label{eq:excision-BC-identification}
i_V^!v^*A\ \simeq\ i^!A.
\end{equation}
Combining \eqref{eq:excision-adj-X}, \eqref{eq:excision-adj-V}, and \eqref{eq:excision-BC-identification} yields a canonical isomorphism
\begin{equation}\label{eq:excision-supports-iso}
H_Z^d(X;A)\xrightarrow{\ \sim\ } H_Z^d(V;v^*A).
\end{equation}
Since the forget-support maps are induced by the counits $i_*i^!\Rightarrow \id$ and $i_{V*}i_V^!\Rightarrow \id$, and Beck--Chevalley is compatible with these counits, the isomorphism \eqref{eq:excision-supports-iso} carries supported refinements of $c$ bijectively to supported refinements of $v^*c$. Thus
$$
\SLoc_Z^d(c)\xrightarrow{\ \sim\ }\SLoc_Z^d(v^*c),
$$
and after adjunction this identifies the localisation torsors
$$
\Loc_Z^{\mathrm{tor},X}(c)
\ =\
\Loc_Z^{\mathrm{tor},V}(v^*c)
$$
under the stated identification of targets.
Finally, if either torsor is a singleton, the equality of torsors forces equality of the unique element, giving
$$
\Loc_Z^{X}(c)=\Loc_Z^{V}(v^*c).
$$
\end{proof}

\end{document}
